\documentclass[11pt]{amsart}
\usepackage[margin=1.15in]{geometry}
\usepackage{amsmath,amssymb,amsthm,mathtools}
\usepackage[T1]{fontenc}
\usepackage{lmodern}
\usepackage{microtype}
\usepackage{enumitem}
\usepackage{alltt}
\usepackage{booktabs}
\usepackage{xcolor}
\usepackage[colorlinks=true,linkcolor=blue!50!black,citecolor=blue!50!black,urlcolor=blue!50!black]{hyperref}

\newtheorem{theorem}{Theorem}[section]
\newtheorem{lemma}[theorem]{Lemma}
\newtheorem{proposition}[theorem]{Proposition}
\newtheorem{corollary}[theorem]{Corollary}
\newtheorem{statement}[theorem]{Statement}
\newtheorem{maintheorem}{Theorem}
\renewcommand{\themaintheorem}{\Alph{maintheorem}}
\newtheorem{mainproposition}[maintheorem]{Proposition}
\theoremstyle{definition}
\newtheorem{definition}[theorem]{Definition}
\newtheorem{question}{Question}
\renewcommand{\thequestion}{\Alph{question}}
\theoremstyle{remark}
\newtheorem{remark}[theorem]{Remark}

\newcommand{\R}{\mathbb{R}}
\newcommand{\N}{\mathbb{N}}
\newcommand{\dimH}{\dim_{\mathrm H}}
\newcommand{\Lm}{\mathcal{L}}
\newcommand{\Dy}{\mathcal{D}}
\newcommand{\G}{\mathcal{G}_q}
% Lean declaration names (underscores allowed)
\newcommand{\lean}[1]{\texttt{\detokenize{#1}}}
\newcommand{\leanref}[1]{\hfill{\footnotesize[Lean: \lean{#1}]}}

\title[Erd\H{o}s similarity for distance sets]{Erd\H{o}s similarity for distance sets:\\
finite patterns, a pinned obstruction, and an open window}
\author{Harald Schilly}
\address{Independent researcher, Austria}
\email{harald.schilly@gmail.com}
\date{October 2026}

\begin{document}

\begin{abstract}
For a compact set $E\subset\R^d$, $d\ge2$, let $\Delta(E)=\{|x-y|:x,y\in E\}$ be its distance set.
We ask whether distance sets of sets of large Hausdorff dimension must contain a nontrivial
affine copy $x+s\{2^{-n}:n\ge1\}$ of the dyadic sequence. This is a distance-set version of
the Erd\H{o}s similarity problem. We combine two recent results announced by OpenAI: the Falconer
distance theorem in all dimensions, and the dyadic case of the Erd\H{o}s similarity conjecture.
Assuming these, we show the following. (A) If $\dimH E>d/2$, then $\Delta(E)$ contains affine
copies of every finite set, for a positive-measure family of translations at every small scale.
(B) There is a compact set $K\subset\R^d$ of nearly full volume in the unit ball whose distance
set from the origin contains no affine copy of the dyadic sequence. The full distance set
$\Delta(K)$ does contain one, and so does the distance set from every other point. (C) A negative
answer to the question forces $\Delta(E)$ to have positive measure but empty interior; with the
classical Mattila--Sj\"olin theorem, the question is therefore open exactly for
$d/2<\dimH E\le(d+1)/2$. Every statement of this note is proved in Lean~4 with Mathlib. The two
OpenAI results enter as explicit hypotheses, and we re-checked OpenAI's own Lean proofs of them
with the Lean kernel, so Theorems~A and~B and Proposition~C are machine-checked down to the
standard axioms. The Mattila--Sj\"olin theorem, which is not formalized, is used only in one
clearly marked remark.
\end{abstract}

\subjclass[2020]{28A80 (primary); 28A75, 42B10, 03B35}
\keywords{Erd\H{o}s similarity problem, distance sets, Falconer conjecture, Hausdorff dimension,
formal verification, Lean}

\maketitle

\section{Introduction}

\subsection{Two problems about large sets}
A set $P\subset\R$ is \emph{measure universal} if every measurable $S\subset\R$ of positive
Lebesgue measure contains a nontrivial affine copy $x+sP$ ($x\in\R$, $s\neq0$). Steinhaus'
theorem~\cite{Steinhaus1920} implies that every finite set is measure universal. Erd\H{o}s
conjectured~\cite{Erdos1974} that no infinite set is. The conjecture is known for slowly
decaying sequences~\cite{Falconer1984} and for many structured sets~\cite{Bourgain1987,Kolountzakis1997}.
The geometric sequences $\{q^n\}$, and above all the dyadic sequence
\[
  \Dy=\{2^{-n}:n\ge1\},
\]
were long described as the main bottleneck~\cite{JungLaiMooroogen2024}. In September and October
2026, OpenAI released a collection of manuscripts produced by an internal AI model. It
includes a proof that $\Dy$ is not measure universal~\cite{OAIDyadic} and an extension to every
ratio $q\in(0,1)$~\cite{OAIGeometric}.

Falconer's distance problem~\cite{Falconer1985} asks for the smallest $\alpha$ such that every
compact $E\subset\R^d$ with $\dimH E>\alpha$ has a distance set
\[
  \Delta(E)=\{|x-y|:x,y\in E\}
\]
of positive Lebesgue measure. Falconer showed that $d/2$ would be optimal and conjectured it.
Partial progress includes~\cite{Wolff1999,Erdogan2005,GIOW2020}. The same OpenAI release contains a
proof of the full conjecture in every dimension $d\ge2$~\cite{OAIFalconer}.
Mattila and Sj\"olin~\cite{MattilaSjolin1999} showed that $\dimH E>(d+1)/2$ even forces
$\Delta(E)$ to contain an interval.

\subsection{The question}
The Falconer theorem says that $\Delta(E)$ is large in measure, while the dyadic theorem says
that being large in measure is not enough to contain $\Dy$. Distance sets, however, have
additional multi-scale structure: $\Delta(E)\supseteq\Delta(E\cap B)$ for every ball $B$. In
particular (Lemma~\ref{lem:scales}), if $\dimH E>d/2$ then $\Delta(E)\cap[0,r]$ has positive
measure for every $r>0$. This motivates:

\begin{question}[Erd\H{o}s similarity for distance sets]\label{q:main}
Let $d\ge2$ and let $E\subset\R^d$ be compact with $\dimH E>d/2$. Must $\Delta(E)$ contain a
nontrivial affine copy of $\Dy$?
\end{question}

We could not find this question, or any study of Erd\H{o}s-type patterns inside distance sets,
in the literature; it is not discussed in the recent survey~\cite{JungLaiMooroogen2024}.

\subsection{Results}
Theorem~\ref{thm:A} assumes the Falconer theorem (Theorem~\ref{thm:inputFalconer}), and
Theorem~\ref{thm:B} assumes the dyadic avoidance theorem (Theorem~\ref{thm:inputDyadic}); both
inputs are stated precisely in Section~\ref{sec:inputs}.

\begin{maintheorem}[Finite patterns]\label{thm:A}
Assume Theorem~\ref{thm:inputFalconer}. Let $d\ge2$, let $E\subset\R^d$ be compact with $\dimH E>d/2$, and let
$F\subset\R$ be finite. There is $\delta>0$ such that for every $|s|<\delta$,
\[
  \Lm^1\bigl(\{x\in\R : x+sF\subseteq\Delta(E)\}\bigr)>0 .
\]
In particular $\Delta(E)$ contains affine copies $x+sF$ with $0<s<\delta'$ and with
$-\delta'<s<0$, for every $\delta'>0$.
\end{maintheorem}

\begin{maintheorem}[A pinned obstruction]\label{thm:B}
Assume Theorem~\ref{thm:inputDyadic}. Let $d\ge1$ and $\eta\in(0,1)$. There is a compact set $K$ in the closed
unit ball of $\R^d$ with $0\in K$ and
\[
  \Lm^d(K)\ \ge\ (1-\eta)^d\,\Lm^d(B(0,1))>0,\qquad\text{hence }\dimH K=d,
\]
such that the pinned distance set $\Delta_0(K)=\{|y|:y\in K\}$ contains no nontrivial affine
copy of $\Dy$, whereas the full distance set $\Delta(K)$ contains one.
If $d\ge2$, then for every $p\in\R^d\setminus\{0\}$ the pinned distance set
$\Delta_p(K)=\{|p-y|:y\in K\}$ contains a nondegenerate interval, hence an affine copy of $\Dy$.
\end{maintheorem}

So pinned distance sets are not universal hosts for $\Dy$, even from a pin inside a set of
full dimension and nearly full volume in the ball. The obstruction, however, sits at one single
point. Theorem~\ref{thm:B} holds verbatim for every ratio $q\in(0,1)$ in place of $1/2$, assuming
the corresponding avoidance statement (Theorem~\ref{thm:Bq}).

\begin{mainproposition}\label{prop:C}
Let $d\ge2$, let $E\subset\R^d$ be compact, and let $q\in(0,1)$.
\begin{enumerate}[label=\textup{(\roman*)},leftmargin=*]
\item If $\Delta(E)$ has nonempty interior, it contains an affine copy of every bounded set
  $P\subseteq[0,M]$ and of $\{q^n:n\ge1\}$.
\item This applies whenever $\Lm^d(E)>0$.
\item Assume Theorem~\ref{thm:inputFalconer}. If $\dimH E>d/2$ and $\Delta(E)$ contains no affine
  copy of $\{q^n:n\ge1\}$, then $\Lm^1(\Delta(E))>0$, $\Delta(E)$ has empty interior, and
  $\Lm^d(E)=0$.
\end{enumerate}
\end{mainproposition}

Combined with the Mattila--Sj\"olin theorem, a counterexample to Question~\ref{q:main} must have
$d/2<\dimH E\le(d+1)/2$ (Remark~\ref{rem:MS}), so the question is open exactly in this window. We
are not aware of a general result that guarantees nonempty interior of $\Delta(E)$ below the
Mattila--Sj\"olin threshold $(d+1)/2$ (see~\cite{MattilaBook2015} for background). Under that reading, a
negative answer to Question~\ref{q:main} would also give new examples of distance sets with
positive measure and empty interior. Section~\ref{sec:open} lists further questions.

\subsection{Formalization and status of the inputs}\label{sec:status}
\emph{Every numbered statement of this note is proved in Lean~4~\cite{Lean4} with
Mathlib~\cite{Mathlib2020}}. The Lean name is given next to each statement and in
Table~\ref{tab:lean}. External results enter only as explicit hypotheses, never as axioms:
\begin{enumerate}[leftmargin=*]
\item the Falconer theorem~\cite{OAIFalconer} (Theorem~\ref{thm:inputFalconer}), used in
  Theorem~\ref{thm:A}, Lemma~\ref{lem:scales}, Proposition~\ref{prop:pairs} and
  Proposition~\ref{prop:C}(iii);
\item the dyadic avoidance theorem~\cite{OAIDyadic} (Theorem~\ref{thm:inputDyadic}), used in
  Theorem~\ref{thm:B};
\item the geometric avoidance statement for a ratio $q\ne1/2$~\cite{OAIGeometric}
  (Statement~\ref{st:geom}), used only in Theorem~\ref{thm:Bq};
\item the Mattila--Sj\"olin theorem~\cite{MattilaSjolin1999} (Statement~\ref{st:MS}), used only
  in Remark~\ref{rem:MS}.
\end{enumerate}
Inputs~(1) and~(2) come from an AI-generated, unrefereed manuscript collection whose authors state
that some results ``could have issues''. OpenAI also provides Lean proofs of both. We rebuilt the
relevant part of that library ourselves and checked with the Lean kernel that these proofs
establish exactly our hypotheses (Section~\ref{sec:lean}). Consequently Theorems~\ref{thm:A}
and~\ref{thm:B}, Proposition~\ref{prop:C} and all lemmas are machine-checked down to Lean's standard
axioms. Input~(3) has no Lean proof, and input~(4) is a classical peer-reviewed theorem that has not
been formalized. A formal proof certifies the statement as formalized; it does not replace
independent mathematical review of~\cite{OAIFalconer,OAIDyadic}.

\section{Notation and inputs}\label{sec:inputs}
$\Lm^d$ is Lebesgue measure on $\R^d$, $|\cdot|$ the Euclidean norm, $B(0,1)$ the open unit
ball, and $\dimH$ Hausdorff dimension. For $E\subset\R^d$ and $p\in\R^d$ we write
$\Delta(E)=\{|x-y|:x,y\in E\}$ and $\Delta_p(E)=\{|p-y|:y\in E\}$. An \emph{affine copy} of
$P\subset\R$ is a set $x+sP$ with $x\in\R$ and $s\ne0$. For $q\in(0,1)$ let
$\G=\{q^n:n\ge1\}$, so $\Dy=\mathcal G_{1/2}$.

\begin{theorem}[Falconer distance theorem~{\cite{OAIFalconer}}]\label{thm:inputFalconer}
Let $d\ge2$ and let $E\subset\R^d$ be compact with $\dimH E>d/2$. Then $\Lm^1(\Delta(E))>0$.
\leanref{FalconerStatement}
\end{theorem}

\begin{theorem}[Dyadic avoidance~{\cite{OAIDyadic}}]\label{thm:inputDyadic}
For every $\eta\in(0,1)$ there is a compact set $A\subseteq[0,1]$ with $\Lm^1(A)>1-\eta$ such that
for every $x\in\R$ and every $s\neq0$ there is $n\ge1$ with $x+s2^{-n}\notin A$.
\leanref{DyadicAvoidanceStatement}
\end{theorem}

\begin{statement}[Geometric avoidance for the ratio $q$]\label{st:geom}
Statement $\mathrm{GA}(q)$: Theorem~\ref{thm:inputDyadic} with $2^{-n}$ replaced by $q^n$. For
$q=1/2$ it is Theorem~\ref{thm:inputDyadic}; for all $q\in(0,1)$ it is claimed
in~\cite{OAIGeometric}, without a Lean proof. \leanref{GeometricAvoidanceStatement}
\end{statement}

\begin{statement}[Mattila--Sj\"olin~\cite{MattilaSjolin1999}]\label{st:MS}
If $d\ge2$ and $E\subset\R^d$ is compact with $\dimH E>(d+1)/2$, then $\Delta(E)$ has nonempty
interior. \leanref{MattilaSjolinStatement}
\end{statement}

\section{Finite patterns}

The following is the standard Steinhaus-type argument (compare the introduction
of~\cite{OAIDyadic}). We include it because Theorem~\ref{thm:A} needs the version with a
positive-measure set of translations at every small scale.

\begin{lemma}\label{lem:steinhaus}
Let $S\subset\R$ be measurable with $0<\Lm^1(S)<\infty$ and let $F\subset\R$ be finite. There is
$\delta>0$ such that $\Lm^1(\{x:x+sF\subseteq S\})>0$ for all $|s|<\delta$.
\leanref{eventually_volume_pattern_pos}
\end{lemma}

\begin{proof}
For $a\in F$ let $T_{a,s}=S-sa$. Translation is continuous in measure:
$\Lm^1(T_{a,s}\,\triangle\,S)\to0$ as $s\to0$. Choose $\delta>0$ with
$\sum_{a\in F}\Lm^1(T_{a,s}\triangle S)<\Lm^1(S)$ for $|s|<\delta$. Since
\[
  S\subseteq\{x: x+sF\subseteq S\}\ \cup\ \bigcup_{a\in F}\,(S\setminus T_{a,s}),
\]
we get $\Lm^1(S)\le\Lm^1(\{x:x+sF\subseteq S\})+\sum_{a\in F}\Lm^1(T_{a,s}\triangle S)$, and the
claim follows.
\end{proof}

\begin{proof}[Proof of Theorem~\ref{thm:A}]
\leanref{finite_pattern_distSet, exists_finite_pattern_distSet}\\
$\Delta(E)$ is the image of the compact set $E\times E$ under the continuous map
$(x,y)\mapsto|x-y|$, so it is compact. In particular it is measurable with finite measure, and
$\Lm^1(\Delta(E))>0$ by Theorem~\ref{thm:inputFalconer}. Lemma~\ref{lem:steinhaus} with $S=\Delta(E)$ gives the first
claim. For the second, apply it with $s=\pm\min(\delta,\delta')/2$; a set of positive measure is
nonempty.
\end{proof}

\begin{proposition}[Planar form]\label{prop:pairs}
Under the assumptions of Theorem~\ref{thm:A}, the set $\{(s,x)\in\R^2:x+sF\subseteq\Delta(E)\}$ has
positive planar Lebesgue measure. \leanref{volume_pattern_pairs_pos}
\end{proposition}

\begin{proof}
The set is closed, hence measurable. By Tonelli's theorem its measure is
$\int\Lm^1(\{x:x+sF\subseteq\Delta(E)\})\,ds$, and the integrand is positive for $|s|<\delta$.
\end{proof}

\begin{lemma}[Distances at every scale]\label{lem:scales}
Assume Theorem~\ref{thm:inputFalconer}. If $d\ge2$ and $E\subset\R^d$ is compact with $\dimH E>d/2$,
then $\Lm^1(\Delta(E)\cap[0,r])>0$ for every $r>0$.
\leanref{volume_distSet_inter_Icc_pos}
\end{lemma}

\begin{proof}
Cover $E$ by finitely many closed balls $\bar B(c_i,r/2)$. Hausdorff dimension of a finite union
is the maximum of the dimensions, so some compact piece $E_i=E\cap\bar B(c_i,r/2)$ has
$\dimH E_i>d/2$. By Theorem~\ref{thm:inputFalconer}, $\Lm^1(\Delta(E_i))>0$, and
$\Delta(E_i)\subseteq\Delta(E)\cap[0,r]$.
\end{proof}

\begin{remark}
Theorem~\ref{thm:A} uses nothing about distance sets beyond compactness and positive measure. Its
point is the contrast with Theorem~\ref{thm:B} and Question~\ref{q:main}: the finite-to-infinite
gap of the Erd\H{o}s similarity problem reappears inside Falconer distance sets.
\end{remark}

\section{The pinned obstruction}

\begin{lemma}[Adding the limit point]\label{lem:zero}
Let $q\in(0,1)$ and let $A\subset\R$ contain no nontrivial affine copy of $\G$. Then neither does
$A\cup\{0\}$. \leanref{not_containsGeomCopy_insert_zero}
\end{lemma}

\begin{proof}
Suppose $x+sq^n\in A\cup\{0\}$ for all $n\ge1$, with $s\ne0$. The points $x+sq^m$, $m\ge0$,
are pairwise distinct. If none of them is $0$, then $x+s\G\subseteq A$, a contradiction. Otherwise
$x+sq^m=0$ for exactly one $m\ge0$. The tail $\{x+sq^{m+n}:n\ge1\}=x+(sq^m)\G$ is then
an affine copy of $\G$ avoiding $0$, hence contained in $A$, again a contradiction.
\end{proof}

\begin{lemma}[Radial sets]\label{lem:radial}
Let $d\ge1$ and let $A\subseteq[0,1]$ be compact with $\Lm^1(A)=c$. Then
$K_A=\{y\in\R^d:|y|\in A\}$ satisfies $\Lm^d(K_A)\ge c^d\,\Lm^d(B(0,1))$.
\leanref{volume_norm_preimage_ge}
\end{lemma}

\begin{proof}
In polar coordinates,
\[
  \Lm^d(K_A)=d\,\Lm^d(B(0,1))\int_A r^{d-1}\,dr .
\]
Since $r\mapsto r^{d-1}$ is nondecreasing,
\begin{align*}
  \int_A r^{d-1}dr-\int_0^c r^{d-1}dr
  &=\int_{A\setminus[0,c]}r^{d-1}dr-\int_{[0,c]\setminus A}r^{d-1}dr\\
  &\ge c^{d-1}\bigl(\Lm^1(A\setminus[0,c])-\Lm^1([0,c]\setminus A)\bigr)=c^{d-1}(\Lm^1(A)-c)=0,
\end{align*}
so $\int_A r^{d-1}dr\ge c^d/d$.
\end{proof}

\begin{lemma}\label{lem:dim}
If $K\subset\R^d$ has $\Lm^d(K)>0$, then $\dimH K=d$. \leanref{dimH_eq_of_volume_pos}
\end{lemma}

\begin{proof}
For the sup-norm on $\R^d$, $d$-dimensional Hausdorff measure equals $\Lm^d$, and the Euclidean
and sup norms are bi-Lipschitz equivalent. Hence $\mathcal H^d(K)>0$ and $\dimH K\ge d$. The
reverse inequality is trivial.
\end{proof}

\begin{lemma}[Steinhaus]\label{lem:interval}
If $K\subset\R^d$ ($d\ge1$) is compact with $\Lm^d(K)>0$, then $[0,\varepsilon)\subseteq\Delta(K)$
for some $\varepsilon>0$. \leanref{exists_Ico_subset_distSet}
\end{lemma}

\begin{proof}
By Steinhaus' theorem $K-K\supseteq B(0,\varepsilon)$ for some $\varepsilon>0$. For
$0\le r<\varepsilon$ write $re_1=x-y$ with $x,y\in K$; then $r=|x-y|\in\Delta(K)$.
\end{proof}

\begin{lemma}[Other pins]\label{lem:otherpins}
Let $d\ge2$, $\rho>0$, $p\in\R^d\setminus\{0\}$, and let $K\subseteq\R^d$ contain the sphere
$\{|y|=\rho\}$. Then $\Delta_p(K)\supseteq[\,||p|-\rho|,\,|p|+\rho\,]$.
\leanref{Icc_subset_pinnedDistSet}
\end{lemma}

\begin{proof}
The sphere is connected because $d\ge2$, and the continuous function $y\mapsto|p-y|$ on it takes
the values $||p|-\rho|$ at $\rho p/|p|$ and $|p|+\rho$ at $-\rho p/|p|$. By the intermediate
value theorem it takes every value in between.
\end{proof}

\begin{theorem}[Theorem~\ref{thm:B} for a general ratio]\label{thm:Bq}
Let $q\in(0,1)$ and assume $\mathrm{GA}(q)$ (Statement~\ref{st:geom}). Then Theorem~\ref{thm:B}
holds with $\Dy$ replaced by $\G$. \leanref{pinned_geometric_obstruction}
\end{theorem}

\begin{proof}
Let $A$ be given by $\mathrm{GA}(q)$ for $\eta$, and put $K=\{y:|y|\in A\cup\{0\}\}$. The set
$A\cup\{0\}\subseteq[0,1]$ is compact, so $K$ is compact and lies in the closed unit ball, and
$0\in K$. Since $K\supseteq K_A$ and $\Lm^1(A)>1-\eta$, Lemma~\ref{lem:radial} gives
$\Lm^d(K)\ge(1-\eta)^d\Lm^d(B(0,1))>0$, and Lemma~\ref{lem:dim} gives $\dimH K=d$. Further,
$\Delta_0(K)=\{|y|:y\in K\}\subseteq A\cup\{0\}$, which contains no affine copy of $\G$ by
Lemma~\ref{lem:zero}.

Any set $P\subseteq\R$ containing a nondegenerate interval $[a,b]$ contains the affine copy
$a+(b-a)\G$. By Lemma~\ref{lem:interval}, $\Delta(K)$ contains an interval. For $d\ge2$ and $p\ne0$,
pick $\rho\in A$ with $\rho>0$ (possible since $\Lm^1(A)>0$); the sphere $\{|y|=\rho\}$ lies in $K$,
so by Lemma~\ref{lem:otherpins} $\Delta_p(K)$ contains an interval of length $2\min(\rho,|p|)>0$.
\end{proof}

\begin{proof}[Proof of Theorem~\ref{thm:B}]
\leanref{pinned_dyadic_obstruction}\\
Theorem~\ref{thm:Bq} with $q=1/2$, where $\mathrm{GA}(1/2)$ is Theorem~\ref{thm:inputDyadic}.
\end{proof}

\section{Where the question is open}

\begin{proof}[Proof of Proposition~\ref{prop:C}]
\leanref{distSet_containsAffineCopy_of_interior, distSet_containsAffineCopy_of_volume_pos, counterexample_properties}\\
(i) A set with nonempty interior contains an interval $[a,b]$ with $a<b$. For
$P\subseteq[0,M]$, $M>0$, we have $a+\frac{b-a}{M}P\subseteq[a,b]$, and $\G\subseteq[0,1]$.
(ii) follows from Lemma~\ref{lem:interval}. (iii) $\Lm^1(\Delta(E))>0$ is
Theorem~\ref{thm:inputFalconer}; nonempty interior or $\Lm^d(E)>0$ would give a copy of $\G$ by (i)
and (ii).
\end{proof}

\begin{remark}[Uses the Mattila--Sj\"olin theorem]\label{rem:MS}
Assume Statement~\ref{st:MS}. If $d\ge2$, $E\subset\R^d$ is compact and $\Delta(E)$ contains no
affine copy of $\G$, then $\dimH E\le(d+1)/2$: otherwise $\Delta(E)$ has nonempty interior and
Proposition~\ref{prop:C}(i) applies. Together with Proposition~\ref{prop:C}(iii), a counterexample to
Question~\ref{q:main} lies in the window $d/2<\dimH E\le(d+1)/2$.
\leanref{window_of_mattilaSjolin}
\end{remark}

\section{Open questions}\label{sec:open}
\begin{enumerate}[leftmargin=*]
\item Question~\ref{q:main} in the window $d/2<\dimH E\le(d+1)/2$. A positive answer would need
  multi-scale structure of distance sets beyond positive measure. A negative answer would need a
  construction controlling all distances of a set of dimension $>d/2$ simultaneously.
\item \emph{All pins.} Is there a compact $E\subset\R^d$ with $\Lm^d(E)>0$ (or merely $\dimH E>d/2$)
  such that $\Delta_p(E)$ contains no affine copy of $\Dy$ for \emph{every} $p\in E$? By
  Lemma~\ref{lem:spheres} below, this asks for a set whose complement contains, around every point
  of $E$, a sphere of radius $r$ for some $r$ in every affine copy of $\Dy$. This suggests a
  connection with spherical maximal function estimates. Pinned Falconer-type
  theorems~\cite{PeresSchlag2000,GIOW2020} give pins with large pinned distance sets, but say
  nothing about such patterns.
\item The same questions for general infinite sets $P$, for instance for every $P$ known to be
  non-universal.
\end{enumerate}

\begin{lemma}\label{lem:spheres}
For $p\in\R^d$ and $E\subseteq\R^d$, the pinned distance set $\Delta_p(E)$ contains no affine copy
of $\G$ if and only if for all $x\in\R$ and $s\ne0$ there is $n\ge1$ such that the sphere of radius
$x+sq^n$ about $p$ is disjoint from $E$. \leanref{not_containsGeomCopy_pinnedDistSet_iff}
\end{lemma}

\begin{proof}
$r\in\Delta_p(E)$ if and only if the sphere of radius $r$ about $p$ meets $E$.
\end{proof}

\section{Formalization}\label{sec:lean}
\begin{sloppypar}
The Lean development accompanying this note is available at
\url{https://github.com/haraldschilly/math-distance-set-similarity}. It is pinned to Lean~4
\texttt{v4.34.1} and Mathlib commit \texttt{d13f23b}, the same versions as OpenAI's library, so the
input statements typecheck identically. The four hypotheses of Section~\ref{sec:status} are Lean
propositions in \texttt{Statements.lean}. \lean{FalconerStatement} and
\lean{DyadicAvoidanceStatement} are copied verbatim from OpenAI's statement files
\texttt{FalconerAllDimensions.lean} and \texttt{DyadicAvoidance.lean}. Table~\ref{tab:lean} lists every
result of this note with its Lean name and its hypotheses. For example, Theorem~\ref{thm:B} reads:
\end{sloppypar}
\begin{alltt}\footnotesize
theorem pinned_dyadic_obstruction (hDy : DyadicAvoidanceStatement) \{d : \(\mathbb{N}\)\}
    (hd : 1 \(\le\) d) \{\(\eta\) : \(\mathbb{R}\)\} (h\(\eta\)0 : 0 < \(\eta\)) (h\(\eta\)1 : \(\eta\) < 1) :
    \(\exists\) K : Set (EuclideanSpace \(\mathbb{R}\) (Fin d)),
      IsCompact K \(\wedge\) K \(\subseteq\) closedBall 0 1 \(\wedge\)
      (0 : EuclideanSpace \(\mathbb{R}\) (Fin d)) \(\in\) K \(\wedge\)
      ENNReal.ofReal ((1 - \(\eta\)) ^ d) * volume (ball 0 1) \(\le\) volume K \(\wedge\)
      0 < volume K \(\wedge\) dimH K = d \(\wedge\)
      \(\neg\) ContainsDyadicCopy (pinnedDistSet 0 K) \(\wedge\) ContainsDyadicCopy (distSet K) \(\wedge\)
      (2 \(\le\) d \(\to\) \(\forall\) p, p \(\ne\) 0 \(\to\) ContainsDyadicCopy (pinnedDistSet p K))
\end{alltt}

\begin{table}[ht]
\centering\footnotesize
\begin{tabular}{@{}lll@{}}
\toprule
Result & Lean name & Hypotheses \\
\midrule
Lemma~\ref{lem:scales} & \lean{volume_distSet_inter_Icc_pos} & Falconer \\
Lemma~\ref{lem:steinhaus} & \lean{eventually_volume_pattern_pos} & none \\
Theorem~\ref{thm:A} & \lean{finite_pattern_distSet}, \lean{exists_finite_pattern_distSet} & Falconer \\
Proposition~\ref{prop:pairs} & \lean{volume_pattern_pairs_pos} & Falconer \\
Lemma~\ref{lem:zero} & \lean{not_containsGeomCopy_insert_zero} & none \\
Lemma~\ref{lem:radial} & \lean{volume_norm_preimage_ge} & none \\
Lemma~\ref{lem:dim} & \lean{dimH_eq_of_volume_pos} & none \\
Lemma~\ref{lem:interval} & \lean{exists_Ico_subset_distSet} & none \\
Lemma~\ref{lem:otherpins} & \lean{Icc_subset_pinnedDistSet} & none \\
Theorem~\ref{thm:B} & \lean{pinned_dyadic_obstruction} & dyadic avoidance \\
Theorem~\ref{thm:Bq} & \lean{pinned_geometric_obstruction} & $\mathrm{GA}(q)$ \\
Proposition~\ref{prop:C}(i),(ii) & \lean{distSet_containsAffineCopy_of_interior}, \lean{..._of_volume_pos} & none \\
Proposition~\ref{prop:C}(iii) & \lean{counterexample_properties} & Falconer \\
Remark~\ref{rem:MS} & \lean{window_of_mattilaSjolin} & Mattila--Sj\"olin \\
Lemma~\ref{lem:spheres} & \lean{not_containsGeomCopy_pinnedDistSet_iff} & none \\
\bottomrule
\end{tabular}
\caption{Results and their Lean formalizations (namespace \lean{DistanceSimilarity}). Every listed
theorem depends only on the axioms \lean{propext}, \lean{Classical.choice}, \lean{Quot.sound}.}
\label{tab:lean}
\end{table}

\begin{sloppypar}
\emph{Re-checking the inputs.} A script in the repository copies, from a checkout of
\texttt{github.com/openai/math} at commit \texttt{adc7f12}, the transitive closure of OpenAI's Lean
modules proving the two inputs: 2459 files, depending only on Mathlib. These are built with the
same toolchain. A further module states
\texttt{theorem : FalconerStatement := OAI.Falconer.falconer\_distance\_conjecture} and the
analogous dyadic statement. It then derives unconditional versions of Lemma~\ref{lem:scales},
Theorem~\ref{thm:A}, Proposition~\ref{prop:pairs}, Theorem~\ref{thm:B} and
Proposition~\ref{prop:C}(iii). The kernel accepts these, and \texttt{\#print axioms} reports only
the three standard axioms. We additionally replayed the entire resulting environment (Mathlib,
OpenAI's modules and ours) in a fresh kernel with Lean's \texttt{leanchecker}. We did not run
OpenAI's \texttt{comparator} setup. A formal proof certifies the statement as formalized. Readers
should still compare the formal statements with the informal ones; to make this easy, the
repository contains a plain-language guide to reading the Lean statements.
\end{sloppypar}

\section*{Acknowledgments and use of AI}
This note builds on results produced by an AI model and released by OpenAI~\cite{OAIFalconer,OAIDyadic}.
It was itself prepared with substantial assistance from Claude (Anthropic), an AI system, which
helped choose the combination studied here, search the literature, draft the text, and write
the Lean formalization. The author checked the mathematics and the formal statements and takes
full responsibility for the content.

\bibliographystyle{amsplain}
\bibliography{refs}

\end{document}
